% Options for packages loaded elsewhere
\PassOptionsToPackage{unicode}{hyperref}
\PassOptionsToPackage{hyphens}{url}
%
\documentclass[
  ignorenonframetext,
  aspectratio=43,
]{beamer}
\usepackage{pgfpages}
\setbeamertemplate{caption}[numbered]
\setbeamertemplate{caption label separator}{: }
\setbeamercolor{caption name}{fg=normal text.fg}
\beamertemplatenavigationsymbolshorizontal
% Prevent slide breaks in the middle of a paragraph
\widowpenalties 1 10000
\raggedbottom
\setbeamertemplate{part page}{
  \centering
  \begin{beamercolorbox}[sep=16pt,center]{part title}
    \usebeamerfont{part title}\insertpart\par
  \end{beamercolorbox}
}
\setbeamertemplate{section page}{
  \centering
  \begin{beamercolorbox}[sep=12pt,center]{part title}
    \usebeamerfont{section title}\insertsection\par
  \end{beamercolorbox}
}
\setbeamertemplate{subsection page}{
  \centering
  \begin{beamercolorbox}[sep=8pt,center]{part title}
    \usebeamerfont{subsection title}\insertsubsection\par
  \end{beamercolorbox}
}
\AtBeginPart{
  \frame{\partpage}
}
\AtBeginSection{
  \ifbibliography
  \else
    \frame{\sectionpage}
  \fi
}
\AtBeginSubsection{
  \frame{\subsectionpage}
}

\usepackage{amsmath,amssymb}
\usepackage{lmodern}
\usepackage{iftex}
\ifPDFTeX
  \usepackage[T1]{fontenc}
  \usepackage[utf8]{inputenc}
  \usepackage{textcomp} % provide euro and other symbols
\else % if luatex or xetex
  \usepackage{unicode-math}
  \defaultfontfeatures{Scale=MatchLowercase}
  \defaultfontfeatures[\rmfamily]{Ligatures=TeX,Scale=1}
\fi
\usetheme[]{default}
\usecolortheme{default}
% Use upquote if available, for straight quotes in verbatim environments
\IfFileExists{upquote.sty}{\usepackage{upquote}}{}
\IfFileExists{microtype.sty}{% use microtype if available
  \usepackage[]{microtype}
  \UseMicrotypeSet[protrusion]{basicmath} % disable protrusion for tt fonts
}{}
\makeatletter
\@ifundefined{KOMAClassName}{% if non-KOMA class
  \IfFileExists{parskip.sty}{%
    \usepackage{parskip}
  }{% else
    \setlength{\parindent}{0pt}
    \setlength{\parskip}{6pt plus 2pt minus 1pt}}
}{% if KOMA class
  \KOMAoptions{parskip=half}}
\makeatother
\usepackage{xcolor}
\newif\ifbibliography
\setlength{\emergencystretch}{3em} % prevent overfull lines
\setcounter{secnumdepth}{-\maxdimen} % remove section numbering


\providecommand{\tightlist}{%
  \setlength{\itemsep}{0pt}\setlength{\parskip}{0pt}}\usepackage{longtable,booktabs,array}
\usepackage{calc} % for calculating minipage widths
\usepackage{caption}
% Make caption package work with longtable
\makeatletter
\def\fnum@table{\tablename~\thetable}
\makeatother
\usepackage{graphicx}
\makeatletter
\def\maxwidth{\ifdim\Gin@nat@width>\linewidth\linewidth\else\Gin@nat@width\fi}
\def\maxheight{\ifdim\Gin@nat@height>\textheight\textheight\else\Gin@nat@height\fi}
\makeatother
% Scale images if necessary, so that they will not overflow the page
% margins by default, and it is still possible to overwrite the defaults
% using explicit options in \includegraphics[width, height, ...]{}
\setkeys{Gin}{width=\maxwidth,height=\maxheight,keepaspectratio}
% Set default figure placement to htbp
\makeatletter
\def\fps@figure{htbp}
\makeatother

\usepackage{ragged2e}
\usepackage{etoolbox}

\apptocmd{\frame}{}{\justifying}{} 

\setbeamertemplate{enumerate items}[default]
\setbeamertemplate{itemize items}[circle]
\setbeamersize{text margin left=2.25em, text margin right=2.25em}

\let\OLDitemize\itemize
\renewcommand\itemize{
    \OLDitemize\addtolength{\itemsep}{3pt plus 1pt minus 1pt}
}
\let\OLDenumerate\enumerate
\renewcommand\enumerate{
    \OLDenumerate\addtolength{\itemsep}{3pt plus 1pt minus 1pt}
}

\usepackage{orcidlink}
\usepackage{amsmath}
\DeclareMathOperator*{\argmax}{argmax}

\makeatletter
\def\@listii{\leftmargin\leftmarginii
              \topsep    1ex
              \parsep    0\p@   \@plus\p@
              \itemsep   \parsep}
\makeatother

\usepackage{lipsum}

\let\olditem\item
\renewcommand\item{\olditem\justifying}

\definecolor{blue}{RGB}{0,114,178}
\definecolor{red}{RGB}{213,94,0}
\definecolor{yellow}{RGB}{240,228,66}
\definecolor{green}{RGB}{0,158,115}

\definecolor{MyBackground}{RGB}{255,253,218}

\setbeamercolor{section in toc}{fg=blue}
\setbeamercolor{subsection in toc}{fg=red}

\setbeamercolor{frametitle}{fg=blue}
\setbeamercolor{title}{fg=blue}

\setbeamertemplate{footline}[frame number]
\setbeamertemplate{navigation symbols}{} 
\setbeamertemplate{itemize items}{-}

\setbeamercolor{itemize item}{fg=blue}
\setbeamercolor{itemize subitem}{fg=blue}
\setbeamercolor{enumerate item}{fg=blue}
\setbeamercolor{enumerate subitem}{fg=blue}
\setbeamercolor{button}{bg=MyBackground,fg=blue}

\newcounter{daggerfootnote}
\newcommand*{\daggerfootnote}[1]{%
    \setcounter{daggerfootnote}{\value{footnote}}%
    \renewcommand*{\thefootnote}{\fnsymbol{footnote}}%
    \footnote[2]{#1}%
    \setcounter{footnote}{\value{daggerfootnote}}%
    \renewcommand*{\thefootnote}{\arabic{footnote}}%
}

\AtBeginSection{
    \begingroup
    \setbeamercolor{background canvas}{bg=yellow}
    \begin{frame}
        \begin{centering}
            \usebeamerfont{section title}\huge\color{black}\insertsection\par
        \end{centering}
    \end{frame}
    \endgroup
}
\makeatletter
\makeatother
\makeatletter
\makeatother
\makeatletter
\@ifpackageloaded{caption}{}{\usepackage{caption}}
\AtBeginDocument{%
\ifdefined\contentsname
  \renewcommand*\contentsname{Table of contents}
\else
  \newcommand\contentsname{Table of contents}
\fi
\ifdefined\listfigurename
  \renewcommand*\listfigurename{List of Figures}
\else
  \newcommand\listfigurename{List of Figures}
\fi
\ifdefined\listtablename
  \renewcommand*\listtablename{List of Tables}
\else
  \newcommand\listtablename{List of Tables}
\fi
\ifdefined\figurename
  \renewcommand*\figurename{Figure}
\else
  \newcommand\figurename{Figure}
\fi
\ifdefined\tablename
  \renewcommand*\tablename{Table}
\else
  \newcommand\tablename{Table}
\fi
}
\@ifpackageloaded{float}{}{\usepackage{float}}
\floatstyle{ruled}
\@ifundefined{c@chapter}{\newfloat{codelisting}{h}{lop}}{\newfloat{codelisting}{h}{lop}[chapter]}
\floatname{codelisting}{Listing}
\newcommand*\listoflistings{\listof{codelisting}{List of Listings}}
\makeatother
\makeatletter
\@ifpackageloaded{caption}{}{\usepackage{caption}}
\@ifpackageloaded{subcaption}{}{\usepackage{subcaption}}
\makeatother
\makeatletter
\@ifpackageloaded{tcolorbox}{}{\usepackage[many]{tcolorbox}}
\makeatother
\makeatletter
\@ifundefined{shadecolor}{\definecolor{shadecolor}{rgb}{.97, .97, .97}}
\makeatother
\makeatletter
\makeatother
\ifLuaTeX
  \usepackage{selnolig}  % disable illegal ligatures
\fi
\IfFileExists{bookmark.sty}{\usepackage{bookmark}}{\usepackage{hyperref}}
\IfFileExists{xurl.sty}{\usepackage{xurl}}{} % add URL line breaks if available
\urlstyle{same} % disable monospaced font for URLs
\hypersetup{
  pdftitle={Seminar 6},
  pdfauthor={Gabriel E. Cabrera Guzmán},
  hidelinks,
  pdfcreator={LaTeX via pandoc}}

\title[Seminar 6]{Seminar
6\daggerfootnote{\texttt{gabriel.cabreraguzman@postgrad.manchester.ac.uk}}}
\subtitle{Statistics 2 (Hypothesis Testing) Solutions}
\author[Gabriel E. Cabrera Guzmán]{
    {Gabriel E. Cabrera
Guzmán\textsuperscript{\normalsize\orcidlink{0000-0000-0000-0000}}} \medskip
}
\date{November 20, 2025}
\institute[]{
    The University of Manchester\\
Alliance Manchester Business School \bigskip
}

\definecolor{quarto-callout-note-color}{RGB}{0,114,178}
\definecolor{quarto-callout-note-color-frame}{RGB}{0,114,178}
\definecolor{quarto-callout-warning-color}{RGB}{213,94,0}
\definecolor{quarto-callout-warning-color-frame}{RGB}{213,94,0}
\begin{document}
\frame{\titlepage}
\ifdefined\Shaded\renewenvironment{Shaded}{\begin{tcolorbox}[interior hidden, breakable, enhanced, sharp corners, boxrule=0pt, frame hidden, borderline west={3pt}{0pt}{shadecolor}]}{\end{tcolorbox}}\fi

\begin{frame}{Exercise 1}
\protect\hypertarget{exercise-1}{}
An international airport advertises that an average of 1,300 flights
land there each day. An aviation analyst suspects that the true average
number of daily landings is higher than this claim. Assuming a
population standard deviation of 720 flights, the analyst monitors
airport activity for 36 days and observes a sample mean of 1,600 flights
per day.

\begin{enumerate}
\item
  Test the claim of the aviation analyst using a 5\% level of
  significance.

  \pause

  Since the population standard deviation is known (\(\sigma = 720\))
  and the sample size is large (\(n = 36\)), we can use the standard
  normal distribution to conduct the hypothesis test.
\end{enumerate}
\end{frame}

\begin{frame}{Exercise 1 (Cont'd)}
\protect\hypertarget{exercise-1-contd}{}
\pause

We first write down the hypotheses. Since the aviation analyst expects
the true average number of daily landings to be higher, the alternative
hypothesis uses a greater-than sign: \vspace{-5pt}

\pause

\[
\begin{aligned}
H_0: \mu &= 1300 \\ 
H_1: \mu &> 1300
\end{aligned}
\]

\pause

Next, we compute the \(z\)-test statistic based on the information
provided. We have \(\mu = 1300\), \(\bar{x} = 1600\), \(n = 36\), and
\(\sigma = 720\):

\pause

\[
z = \frac{\bar{x} - \mu}{\sigma / \sqrt{n}}
= \frac{1600 - 1300}{720 / \sqrt{36}}
= \frac{300}{120}
= 2.5.
\]
\end{frame}

\begin{frame}{Exercise 1 (Cont'd)}
\protect\hypertarget{exercise-1-contd-1}{}
\pause

The significance level is \(\alpha = 5\%\). To find the critical value
for rejecting the null at this level, we look for the value \(z\) such
that the probability to its right is \(0.05\). From the standard normal
table, \vspace{-4pt}

\pause

\[
\Pr(Z < 1.64) = 0.95.
\]

\pause

Thus, the critical value for a right-tailed test at the 5\% level is
\(z_{\alpha} = 1.64\). Since the test statistic (2.5) is greater than
the critical value (1.64), we reject the null hypothesis at the 5\%
significance level.

\pause

Furthermore, the \(p\)-value is the area to the right of the observed
\(z\)-statistic:
\end{frame}

\begin{frame}{Exercise 1 (Cont'd)}
\protect\hypertarget{exercise-1-contd-2}{}
\pause

\[
\begin{aligned}
p &= \Pr(Z > 2.5) \\
  &= 1 - \Pr(Z < 2.5) \\
  &= 1 - 0.9938 \\
  &= 0.0062.
\end{aligned}
\]

\pause

Since the \(p\)-value (\(0.0062\)) is less than \(\alpha = 0.05\), we
again reject the null hypothesis at the 5\% level of significance.
\end{frame}

\begin{frame}{Exercise 1 (Cont'd)}
\protect\hypertarget{exercise-1-contd-3}{}
\begin{enumerate}
\setcounter{enumi}{1}
\item
  Suggest a potential source of a Type I error.

  \pause

  Sample selection bias. The days could be chosen from peak seasons like
  summer.
\end{enumerate}
\end{frame}

\begin{frame}{Exercise 2}
\protect\hypertarget{exercise-2}{}
An insurance firm claims that the monthly number of approved claims is
normally distributed with a mean of 120. An analyst believes that the
true average number of approved claims is less than this and collects
data for the past 25 months. The analyst finds that the sample mean
number of approved claims per month is 114, with a sample standard
deviation of 15.

\begin{enumerate}
\item
  Test the analyst's claim using a 5\% level of significance. You do not
  need to calculate the p-value.

  \pause

  Since the sample standard deviation is known (\(s = 15\)) and the
  population distribution is assumed to be normally distributed, we use
  the Student's \(t\) distribution to conduct the hypothesis test.
\end{enumerate}
\end{frame}

\begin{frame}{Exercise 2 (Cont'd)}
\protect\hypertarget{exercise-2-contd}{}
\pause

We first write down the hypotheses. Since the analyst believes the true
average number of approved claims is lower, the alternative hypothesis
uses a less-than sign:\vspace{-4pt}

\pause

\[
\begin{aligned}
H_0: \mu &= 120 \\
H_1: \mu &< 120
\end{aligned}
\]

\pause

Next, we compute the \(t\)-test statistic based on the information
provided. We have \(\mu = 120\), \(\bar{x} = 114\), \(n = 25\), and
\(s = 15\):

\pause

\[
t = \frac{\bar{x} - \mu}{s / \sqrt{n}}
= \frac{114 - 120}{15 / \sqrt{25}}
= \frac{-6}{3}
= -2.
\]
\end{frame}

\begin{frame}{Exercise 2 (Cont'd)}
\protect\hypertarget{exercise-2-contd-1}{}
\pause

The significance level is \(\alpha = 5\%\). To find the critical value
for rejecting the null at this significance level, we look for the
probability of \(0.05\) in the Student's \(t\) distribution table with
degrees of freedom \vspace{-4pt}

\pause

\[
df = n - 1 = 25 - 1 = 24.
\]

\pause

For a left-tailed test, we look for the value \(t\) such that the
probability to its left is \(0.05\). From the table, we find that
\vspace{-4pt}

\pause

\[
\Pr(T < -1.711) = 0.05.
\]

\pause

Therefore, the critical value for a left-tailed test at the 5\%
significance level is \(-1.711\). Since the \(t\)-test statistic
(\(-2\)) is less than the critical value (\(-1.711\)), we reject the
null hypothesis at the 5\% level of significance.
\end{frame}

\begin{frame}{Exercise 2 (Cont'd)}
\protect\hypertarget{exercise-2-contd-2}{}
\begin{enumerate}
\setcounter{enumi}{1}
\item
  Test the analyst's claim using a 1\% level of significance. You do not
  need to calculate the p-value.

  \pause

  The significance level is \(\alpha = 1\%\). To find the critical value
  for rejecting the null at this significance level, we look for the
  probability of \(0.01\) in the Student's \(t\) distribution table with
  degrees of freedom \vspace{-4pt}

  \pause

  \[
   df = n - 1 = 25 - 1 = 24.
   \]
\end{enumerate}
\end{frame}

\begin{frame}{Exercise 2 (Cont'd)}
\protect\hypertarget{exercise-2-contd-3}{}
\pause

For a left-tailed test, we look for the value \(t\) such that the
probability to its left is \(0.01\). From the table, we find that
\vspace{-4pt}

\pause

\[
\Pr(T < -2.492) = 0.01.
\]

\pause

Therefore, the critical value for a left-tailed test at the 1\%
significance level is \(-2.492\). Since the \(t\)-test statistic
(\(-2\)) is greater than the critical value (\(-2.492\)), we fail to
reject the null hypothesis at the 1\% level of significance.
\end{frame}

\begin{frame}{Exercise 2 (Cont'd)}
\protect\hypertarget{exercise-2-contd-4}{}
\begin{enumerate}
\setcounter{enumi}{2}
\item
  Explain differences (if any) in the results in (a) and (b).

  \pause

  We reject the null hypothesis in part (a) but not in part (b). This
  difference arises because the significance level is lower in part (b).
  A 1\% significance level requires a more extreme test statistic to
  reject the null compared to a 5\% significance level. Since the
  observed test statistic is not extreme enough to exceed the stricter
  1\% critical threshold, we fail to reject the null at 1\% even though
  we were able to reject it at 5\%.
\end{enumerate}
\end{frame}

\begin{frame}{Exercise 3}
\protect\hypertarget{exercise-3}{}
The scores on a first-year undergraduate statistics exam are assumed to
be normally distributed with a population mean of 60 and a population
standard deviation of 18. A sample of 81 students has a sample mean
score of 64.

\begin{enumerate}
\item
  Test the claim that the average exam score is less than 60 using a 1\%
  level of significance.

  \pause

  Since the population standard deviation is known (\(\sigma = 18\)) and
  the sample size is large (\(n = 81\)), we can use the standard normal
  distribution to conduct hypothesis testing. As only the direction of
  the test changes from parts (a)--(c), we first compute the \(z\)-test
  statistic based on the information. We have \(\mu = 60\),
  \(\bar{x} = 64\), \(n = 81\), and \(\sigma = 18\): \vspace{-4pt}

  \pause

  \[
   z = \frac{\bar{x} - \mu}{\sigma / \sqrt{n}}
   = \frac{64 - 60}{18 / \sqrt{81}}
   = \frac{4}{2}
   = 2.
   \]
\end{enumerate}
\end{frame}

\begin{frame}{Exercise 3 (Cont'd)}
\protect\hypertarget{exercise-3-contd}{}
\pause

Since claim is the average exam score is less than 60, the alternative
hypothesis would contain the less than sign. We have: \vspace{-4pt}

\pause

\[
\begin{aligned}
H_0:\mu &= 60 \\
H_1:\mu &< 60
\end{aligned}
\]

\pause

The significance level is \(\alpha=1\%\). To find the critical value for
rejecting the null at this significance level, we look for the
probability of \(0.01\) in the standard normal distribution table. For a
left-tailed test, we look for the value \(z\) such that the probability
to its left is \(0.01\). From the table, we find that
\end{frame}

\begin{frame}{Exercise 3 (Cont'd)}
\protect\hypertarget{exercise-3-contd-1}{}
\pause

\[
\Pr(Z<-2.33)=0.01
\]

\pause

Therefore, the critical value for a left-tailed test at the \(1\%\)
significance level is \(-2.33\). Since the \(z\)-test statistic (2) is
greater than the critical value (\(-2.33\)), we fail to reject the null
hypothesis at the \(1\%\) level of significance.

\pause

Further, the p-value for this hypothesis test is the area to the left of
the observed z-statistic under the standard normal distribution. That
is, \vspace{-4pt}

\pause

\[
\Pr(Z<2)=0.9772
\]
\end{frame}

\begin{frame}{Exercise 3 (Cont'd)}
\protect\hypertarget{exercise-3-contd-2}{}
\pause

Since the p-value (\(0.9772\)) is greater than the significance level of
\(0.01\), we fail to reject the null hypothesis at the \(1\%\) level of
significance.
\end{frame}

\begin{frame}{Exercise 3 (Cont'd)}
\protect\hypertarget{exercise-3-contd-3}{}
\begin{enumerate}
\setcounter{enumi}{1}
\item
  Test the claim that the average exam score is greater than 60 using a
  1\% level of significance.

  \pause

  Since claim is the average exam score is greater than 60, the
  alternative hypothesis would contain the less than sign. We have:
  \vspace{-17pt}

  \pause

  \[
   \begin{aligned}
   H_0:\mu &= 60 \\ 
   H_1:\mu &> 60
   \end{aligned}
   \]

  \pause

  The significance level is \(\alpha=1\%\). To find the critical value
  for rejecting the null at this significance level, we look for the
  probability of \(0.01\) in the standard normal distribution table. For
  a right-tailed test, we look for the value \(z\) such that the
  probability to its right is \(0.01\). From the table, we find that
\end{enumerate}
\end{frame}

\begin{frame}{Exercise 3 (Cont'd)}
\protect\hypertarget{exercise-3-contd-4}{}
\pause

\[
\Pr(Z>2.33)=0.01
\]

\pause

Therefore, the critical value for a right-tailed test at the \(1\%\)
significance level is \(2.33\). Since the \(z\)-test statistic (2) is
less than the critical value (\(2.33\)), we fail to reject the null
hypothesis at the \(1\%\) level of significance.

\pause

Further, the p-value for this hypothesis test is the area to the right
of the observed z-statistic under the standard normal distribution. That
is, \vspace{-10pt}

\pause

\[
\Pr(Z>2)=1-\Pr(Z<2)=1-0.9772=0.023.
\]

\pause

Since the p-value (0.023) is greater than the significance level of
0.01, we fail to reject the null hypothesis at the 1\% level of
significance.
\end{frame}

\begin{frame}{Exercise 3 (Cont'd)}
\protect\hypertarget{exercise-3-contd-5}{}
\begin{enumerate}
\setcounter{enumi}{2}
\item
  Test the claim that the average exam score is different from 60 using
  a 1\% level of significance.

  \pause

  Since claim is the average exam score is different than 60, the
  alternative hypothesis would contain the not equal to sign. We have:
  \vspace{-5pt}

  \pause

  \[
   \begin{aligned}
   H_0:\mu &=60 \\ 
   H_1:\mu &\neq 60
   \end{aligned}
   \]

  \pause

  The significance level is \(\alpha=1\%\). To find the critical value
  for rejecting the null at this significance level, we look for the
  probability of \(0.01 \div 2 = 0.005\) in the standard normal
  distribution table. For a two-tailed test, we look for the value \(z\)
  such that the probability to both its left and right is \(0.005\).
  From the table, we find that
\end{enumerate}
\end{frame}

\begin{frame}{Exercise 3 (Cont'd)}
\protect\hypertarget{exercise-3-contd-6}{}
\pause

\[
\Pr(Z>2.58)=0.005.
\] \vspace{-15pt}

\pause

Therefore, the critical value for a two-tailed test at the \(1\%\)
significance level is \(2.58\). Since the \(z\)-test statistic (2) is
less than the critical value (\(2.58\)), we fail to reject the null
hypothesis at the \(1\%\) level of significance.

\pause

Further, the p-value for this two-tailed hypothesis test is the
probability of observing a test statistic as extreme or more extreme
than the absolute value of the observed z-statistic, in either
direction, under the standard normal distribution. That is,
\vspace{-5pt}

\pause

\[
P(Z > |2|) = 2P(Z > 2) = 2 \times 0.0228 = 0.0456.
\]

\pause

Since the p-value (0.0456) is greater than the significance level of
0.01, we fail to reject the null hypothesis at the 1\% level of
significance.
\end{frame}

\begin{frame}{Exercise 3 (Cont'd)}
\protect\hypertarget{exercise-3-contd-7}{}
\begin{enumerate}
\setcounter{enumi}{3}
\item
  Suppose we change the level of significance level to 5\%. In which of
  the above scenarios can we reject the null?

  When we change the significance level to 5\%, we can reject the null
  for the hypothesis tests in (b) and (c) but not (a).

  \begin{longtable}[]{@{}lll@{}}
  \toprule()
  Test & p-value & Decision \\
  \midrule()
  \endhead
  (a) & 0.977 & Fail to Reject \\
  (b) & 0.023 & Reject \\
  (c) & 0.046 & Reject \\
  \bottomrule()
  \end{longtable}
\end{enumerate}
\end{frame}

\begin{frame}{Exercise 4}
\protect\hypertarget{exercise-4}{}
A manufacturer of rechargeable batteries advertises that the average
lifetime of its batteries is 20 hours. An independent consumer group
wants to investigate this claim. They collect a sample of 100 batteries
and find that the average lifetime in the sample is 19.2 hours. The
population standard deviation is known to be 3 hours.

\begin{enumerate}
\item
  Construct a 99\% confidence interval for the true average battery
  lifetime.

  \pause

  Since the population standard deviation is known and the sample size
  is large. We can use the standard normal distribution to construct the
  confidence interval. To find the confidence interval, we first need to
  find \(Z_{\alpha/2}\). We want to construct a 99\% confidence
  interval. Therefore, the confidence level is \(1-\alpha = 0.99\), and
  \(\alpha = 0.01\). Considering that we have a two-sided confidence
  interval, the error term in each tail must be \(\alpha/2 = 0.005\).
  Thus, \(\Pr(Z < \alpha) = 0.995\).
\end{enumerate}
\end{frame}

\begin{frame}{Exercise 4 (Cont'd)}
\protect\hypertarget{exercise-4-contd}{}
\pause

To find suitable \(Z_{\alpha/2} = \alpha\), search for the probability
of \(0.995\) among the numbers (probabilities) in the standard normal
distribution table. We find that \vspace{-10pt}

\pause

\[
\Pr(Z < 2.576) = 0.995.
\]

Therefore, we have \(Z_{\alpha/2} = 2.576\).
\end{frame}

\begin{frame}{Exercise 4 (Cont'd)}
\protect\hypertarget{exercise-4-contd-1}{}
\pause

Construct the confidence interval and replace \(\bar{x}=19.2\),
\(Z_{\alpha/2}=2.576\), \(n=100\), and \(\sigma=3\) as follows:

\pause

\[
\bar{x} - Z_{\alpha/2} \times \frac{\sigma}{\sqrt{n}}
< \mu <
\bar{x} + Z_{\alpha/2} \times \frac{\sigma}{\sqrt{n}}
\]

\pause

\[
19.2 - 2.576 \times \frac{3}{\sqrt{100}}
< \mu <
19.2 + 2.576 \times \frac{3}{\sqrt{100}}
\]

\pause

\[
18.43 < \mu < 19.97
\]
\end{frame}

\begin{frame}{Exercise 4 (Cont'd)}
\protect\hypertarget{exercise-4-contd-2}{}
\begin{enumerate}
\setcounter{enumi}{1}
\item
  Test the manufacturer's claim using a 1\% level of significance.

  \pause

  Since the 99\% confidence interval does not contain the population
  mean (20) claimed by the manufacturer, we can reject the null
  hypothesis that the population mean is 20 at a 1\% level of
  significance.
\end{enumerate}
\end{frame}

\begin{frame}{Exercise 5}
\protect\hypertarget{exercise-5}{}
A researcher is examining whether the average waiting time at a service
centre has changed from its long-standing average of 15 minutes. A
random sample of 25 customers is selected. The sample shows an average
waiting time of 16.2 minutes.

\pause

\begin{enumerate}
\item
  Assume the population standard deviation is known to be 3 minutes.
  Test whether the true average waiting time differs from 15 minutes
  using a 5\% level of significance.

  \pause

  We first write down the hypothesis, since the claim is that the true
  average waiting time is different, the alternative hypothesis would
  contain the not equal to sign. We have: \vspace{-5pt}

  \pause

  \[
   \begin{aligned}
   H_0:\mu & =15 \\ 
   H_1:\mu & \neq 15
   \end{aligned}
   \]
\end{enumerate}
\end{frame}

\begin{frame}{Exercise 5 (Cont'd)}
\protect\hypertarget{exercise-5-contd}{}
\pause

The population standard deviation is known (\(\sigma = 3\)) and the
sample size is only \(n=25\), we must assume the population distribution
is normal so that we can use the standard normal distribution to conduct
hypothesis testing.

\pause

We compute the z-test statistic based on the information. We have
\(\mu = 15\), \(\bar{x} = 16.2\), \(n = 25\), and \(\sigma = 3\)

\pause

\[
\frac{\bar{x}-\mu}{\sigma/\sqrt{n}}
=
\frac{16.2 - 15}{3/\sqrt{25}}
=
\frac{1.2}{0.6}
=
2
\]
\end{frame}

\begin{frame}{Exercise 5 (Cont'd)}
\protect\hypertarget{exercise-5-contd-1}{}
\pause

The significance level is \(\alpha = 5\%\). To find the critical value
for rejecting the null at this significance level, we look for the
probability of \(0.05 \div 2 = 0.025\) in the standard normal
distribution table. For a two-tailed test, we look for the value \(z\)
such that the probability to both its left and right is \(0.025\). From
the table, we find that \vspace{-5pt}

\pause

\[
\Pr(Z > 1.96) = 0.025.
\]

\pause

Therefore, the critical value for a two-tailed test at the \(5\%\)
significance level is \(1.96\). Since the \(z\)-test statistic (2) is
greater than the critical value (\(1.96\)), we reject the null
hypothesis at the \(5\%\) level of significance.
\end{frame}

\begin{frame}{Exercise 5 (Cont'd)}
\protect\hypertarget{exercise-5-contd-2}{}
\pause

Further, the p-value for this two-tailed hypothesis test is the
probability of observing a test statistic as extreme or more extreme
than the absolute value of the observed z-statistic, in either
direction, under the standard normal distribution. That is,
\vspace{-5pt}

\pause

\[
P(Z > |2|) = 2P(Z > 2) = 2 \times 0.0228 = 0.0456.
\]

\pause

Since the p-value (0.0456) is less than the significance level of 0.05,
we reject the null hypothesis at the 5\% level of significance.
\end{frame}

\begin{frame}{Exercise 5 (Cont'd)}
\protect\hypertarget{exercise-5-contd-3}{}
\begin{enumerate}
\setcounter{enumi}{1}
\item
  Assume the sample standard deviation is known to be 3 minutes. Test
  whether the true average waiting time differs from 15 minutes using a
  5\% level of significance.

  \pause

  Only the sample standard deviation is known (\(s=3\)) and the sample
  size is only \(n=25\), we must assume the population distribution is
  normal so that we can use the Student's \(t\) distribution to conduct
  hypothesis testing.

  \pause

  We compute the t-test statistic based on the information. We have
  \(\mu = 15\), \(\bar{x} = 16.2\), \(n = 25\), and \(s = 3\)

  \pause

  \[
   \frac{\bar{x}-\mu}{s/\sqrt{n}}
   =
   \frac{16.2 - 15}{3/\sqrt{25}}
   =
   \frac{1.2}{0.6}
   =
   2
   \]
\end{enumerate}
\end{frame}

\begin{frame}{Exercise 5 (Cont'd)}
\protect\hypertarget{exercise-5-contd-4}{}
\pause

The significance level is \(\alpha = 5\%\). To find the critical value
for rejecting the null at this significance level, we look for the
probability of \(0.05 \div 2 = 0.025\) in the Student's \(t\)
distribution table with a degree of freedom
\(df = n - 1 = 25 - 1 = 24\). For a two-tailed test, we look for the
value \(t\) such that the probability to both its left and right is
\(0.025\). From the table, we find that \vspace{-5pt}

\pause

\[
\Pr(T > 2.064) = 0.025.
\]

\pause

Therefore, the critical value for a two-tailed test at the \(5\%\)
significance level is \(2.064\). Since the \(t\)-test statistic (2) is
less than the critical value (\(2.064\)), we fail to reject the null
hypothesis at the \(5\%\) level of significance.
\end{frame}

\begin{frame}{Exercise 5 (Cont'd)}
\protect\hypertarget{exercise-5-contd-5}{}
\begin{enumerate}
\setcounter{enumi}{2}
\item
  Explain differences (if any) in the results in (a) and (b).

  \pause

  When using a two-tailed Z-test at the 5\% significance level, the
  critical values are ±1.96. Since the test statistic of 2 exceeds this
  cutoff, we reject the null hypothesis. However, when using a t-test
  with the same test statistic and the same 5\% two-tailed significance
  level, the critical values depend on the degrees of freedom. With a
  smaller sample size, the t-distribution has heavier tails, so the
  critical values are larger in magnitude (for example, around ±2.064
  when df = 24). Because the test statistic of 2 does not exceed these
  larger t critical values, we fail to reject the null hypothesis. Thus,
  the same test statistic can lead to different conclusions because the
  t-distribution is wider than the standard normal distribution when
  sample sizes are small.
\end{enumerate}
\end{frame}



\end{document}
